\section{The exact constant in a change of theta lattice}
\label{tc-section}

The remaining analytic normalization in Section~\ref{lm-section}
can be calculated explicitly.  We give the calculation, including
the finite quadratic Gauss sum on which its eighth power depends.
All square roots of positive real numbers in this section mean
the positive square root.  This convention is part of the
normalization.

Let $V$ be a complex vector space of dimension $r$, let $H$ be a
positive definite Hermitian form, linear in its first variable,
and put $E=\operatorname{Im}H$.  A full lattice $\Lambda$ is
self-dual if
\[
 \Lambda=\Lambda^*
   :=\{x\in V:E(x,\Lambda)\subset\mathbb Z\}.
 \tag{TC1}
\]
A sign semicharacter on $\Lambda$ is a function
$\alpha:\Lambda\to\{1,-1\}$ satisfying
\[
 \alpha(x+y)=\alpha(x)\alpha(y)e^{\pi iE(x,y)}.
 \tag{TC2}
\]
For several lattices we require their semicharacters to be
restrictions of one function, so they agree on intersections.
Write
\[
 U_xF(z)=e^{-\pi H(z,x)-\pi H(x,x)/2}F(z+x),
 \qquad W_x=\alpha(x)U_x.
 \tag{TC3}
\]
The theta space $\mathcal T(\Lambda,\alpha)$ consists of entire
functions $F$ satisfying $W_xF=F$ for $x\in\Lambda$.
The basic identities are
\[
 U_xU_y=e^{-\pi iE(x,y)}U_{x+y},\qquad
 U_xU_y=e^{-2\pi iE(x,y)}U_yU_x.
 \tag{TC4}
\]
They follow by expanding the three exponential factors and using
$H(y,x)=\overline{H(x,y)}$.  On any one of the lattices in
(TC2), the operators $W_x$ form a representation of its additive
group.  Also $\alpha(0)=1$ and $\alpha(-x)=\alpha(x)$.

\subsection{Self-dual lattices and their theta spaces}

Choose Lebesgue measure $dz$ by using any $H$-orthonormal complex
basis.  Thus
\[
                  \int_V e^{-\pi H(z,z)}\,dz=1.
 \tag{TC5}
\]
This measure does not depend on the chosen orthonormal basis.

\begin{lemma}\label{tc-covolume}
Every self-dual lattice has covolume one for this measure.
For every sign semicharacter on such a lattice,
$\dim_{\mathbb C}\mathcal T(\Lambda,\alpha)=1$.
\end{lemma}
\begin{proof}
In real orthonormal coordinates the matrix of $E$ has determinant
one.  If $B$ is a real basis matrix for $\Lambda$, its integral
alternating Gram matrix is $B^{\mathsf T}JB$.
Self-duality says that this integer matrix is unimodular.
Therefore $(\det B)^2=1$, which proves the covolume assertion.

For completeness, the dimension assertion has an elementary
Fourier proof.  A unimodular alternating integer form has a
symplectic basis $e_1,\ldots,e_r,f_1,\ldots,f_r$ with
$E(e_i,f_j)=-\delta_{ij}$.
Indeed, a primitive vector $e_1$ has a partner $f_1$ with
$E(e_1,f_1)=-1$ by unimodularity.  Subtracting integer multiples
of these two vectors splits off their orthogonal complement;
induction supplies the basis.
The $e_i$ are a complex basis of $V$: if
$\sum a_i e_i=i\sum b_i e_i$ with $a_i,b_i\in\mathbb R$,
pairing the two sides with $\sum b_i e_i$ under $E$ gives
$H(\sum b_i e_i,\sum b_i e_i)=0$.  Positivity then gives
$a_i=b_i=0$.

Use this complex basis, and write
$H(z,w)=z^{\mathsf T}G\overline w$ and
$(f_1,\ldots,f_r)=(e_1,\ldots,e_r)Z$.
The matrix $G$ is real symmetric and positive definite, since
$E(e_i,e_j)=0$.  The equations
$E(e_i,f_j)=-\delta_{ij}$ and $E(f_i,f_j)=0$ give
\[
             \operatorname{Im}Z=G^{-1}>0,\qquad Z=Z^{\mathsf T}.
 \tag{TC6}
\]
Choose $a_i,b_i\in\{0,\tfrac12\}$ with
$\alpha(e_i)=e^{2\pi ia_i}$ and
$\alpha(f_i)=e^{2\pi ib_i}$.
For $F\in\mathcal T(\Lambda,\alpha)$ put
$f(z)=e^{-\pi z^{\mathsf T}Gz/2}F(z)$.
The transformation equations become
\[
 \begin{split}
 f(z+u)&=e^{2\pi ia^{\mathsf T}u}f(z),\\
 f(z+Zv)&=e^{2\pi ib^{\mathsf T}v}
        e^{-2\pi iv^{\mathsf T}z-\pi iv^{\mathsf T}Zv}f(z)
                 \qquad(u,v\in\mathbb Z^r).
 \end{split}
 \tag{TC7}
\]
The first equation gives a Fourier expansion
$f(z)=\sum_{n\in\mathbb Z^r}c_n
 e^{2\pi i(n+a)^{\mathsf T}z}$.
Here the coefficients are obtained by integration on a real
unit cube after multiplication by $e^{-2\pi ia^{\mathsf T}z}$.
Cauchy's theorem on the coordinate rectangles makes the
appropriately weighted coefficient independent of the imaginary
part of $z$.  Shifting these rectangles in either direction
bounds the coefficients exponentially on every smaller closed
strip.  The resulting Fourier series therefore converges normally
on compact subsets and represents $f$.
Substitution in the second equation of (TC7) gives
\[
 c_{n+v}=c_n
  \exp\!\left(2\pi i(n+a)^{\mathsf T}Zv+
              \pi iv^{\mathsf T}Zv-2\pi ib^{\mathsf T}v\right).
 \tag{TC8}
\]
All coefficients are determined by $c_0$.
Conversely,
\[
 e^{\pi z^{\mathsf T}Gz/2}
 \sum_{n\in\mathbb Z^r}
  \exp\!\left(\pi i(n+a)^{\mathsf T}Z(n+a)
                 +2\pi i(n+a)^{\mathsf T}(z-b)\right)
 \tag{TC9}
\]
converges normally, because $\operatorname{Im}Z>0$, and
satisfies (TC7).  Its Fourier coefficients are nonzero, so it is
not the zero function.  This proves the dimension assertion.
The case $r=0$ consists simply of the constant functions.
\end{proof}

\subsection{Normalized finite averaging is unitary}

Suppose $\Lambda_i$ are commensurable self-dual lattices carrying
one common sign semicharacter.  Put
\[
 d_{ij}=[\Lambda_j:\Lambda_i\cap\Lambda_j]
        =[\Lambda_i:\Lambda_i\cap\Lambda_j],\qquad
 I_{ij}F=d_{ij}^{-1/2}
     \sum_{x\in\Lambda_j/(\Lambda_i\cap\Lambda_j)}W_xF.
 \tag{TC10}
\]
The equality of indices follows from Lemma~\ref{tc-covolume}.
On functions fixed by $\Lambda_i$ the summands are independent
of the choice of representatives.

\begin{lemma}\label{tc-unitary-intertwiner}
The operator $I_{ij}$ is a unitary isomorphism from the
$\Lambda_i$ theta space to the $\Lambda_j$ theta space, with
inverse $I_{ji}$.  The same inverse and norm assertions hold on
continuous functions with the corresponding transformation laws.
\end{lemma}
\begin{proof}
Let $\Gamma$ be the intersection of all lattices under
consideration.  For continuous functions satisfying the
$\Gamma$ transformation law define
\[
 \langle F,G\rangle_\Gamma
  =\frac1{\operatorname{vol}(V/\Gamma)}
    \int_{V/\Gamma}F(z)\overline{G(z)}
                        e^{-\pi H(z,z)}\,dz.
 \tag{TC11}
\]
The integrand is $\Gamma$-periodic.  Its norm is positive on
nonzero continuous functions.  For $x\in\Gamma^*$ the operator
$U_x$ preserves this space and is unitary, since
\[
 |U_xF(z)|^2e^{-\pi H(z,z)}
       =|F(z+x)|^2e^{-\pi H(z+x,z+x)}.
 \tag{TC12}
\]
All $\Lambda_i$ lie in $\Gamma^*$.
Consequently
\[
       P_i=\frac1{[\Lambda_i:\Gamma]}
                   \sum_{x\in\Lambda_i/\Gamma}W_x
 \tag{TC13}
\]
is an orthogonal projection onto the functions fixed by
$\Lambda_i$: it is the average of a finite unitary group
representation.

For $x\in\Gamma^*$, the commutator in (TC4) and finite
character orthogonality show that $P_iU_xP_i=0$ unless
$x\in\Lambda_i^*=\Lambda_i$.
Explicitly, the character being averaged is
$u\mapsto e^{-2\pi iE(u,x)}$ on $\Lambda_i/\Gamma$;
a nontrivial character has zero sum because translating its
sum by an element on which it is not one multiplies that sum
by a nontrivial scalar.
If $x\in\Lambda_i$, then $U_xP_i=\alpha(x)P_i$.
For $x$ also belonging to another specified lattice, the
common semicharacter therefore gives $W_xP_i=P_i$.
Thus
\[
 P_iP_jP_i=
 \frac{[(\Lambda_i\cap\Lambda_j):\Gamma]}
      {[\Lambda_j:\Gamma]}P_i=d_{ij}^{-1}P_i.
 \tag{TC14}
\]
On the image of $P_i$, regrouping the sum (TC13) gives
$I_{ij}=\sqrt{d_{ij}}P_j$.
Equation (TC14) gives $I_{ji}I_{ij}=1$ and
\[
 \|I_{ij}F\|^2=d_{ij}\langle F,P_jF\rangle
             =\|F\|^2.
\]
The reverse composition is the identity as well.
Finite translations preserve holomorphy, so the assertion
restricts to theta spaces.  For a function satisfying a
$\Lambda_i$ law, its normalized norm in (TC11) is independent
of the smaller common lattice $\Gamma$, by subdivision of
fundamental domains.
\end{proof}

\subsection{An eighth-power formula for finite Gauss sums}

A homogeneous quadratic function on a finite abelian group $M$
is a function $q:M\to\mathbb C^\times$ such that
$q(nx)=q(x)^{n^2}$ for integers $n$ and
\[
                b(x,y)=\frac{q(x+y)}{q(x)q(y)}
 \tag{TC15}
\]
is a bicharacter.  Its radical is
$R=\{x:b(x,M)=1\}$.
All values of $q$ are roots of unity, since
$q(x)^{|M|^2}=q(|M|x)=1$.

\begin{lemma}\label{tc-four-squares-congruence}
For every positive integer $N$ there are integers
$a,b,c,d$ with $a^2+b^2+c^2+d^2\equiv-1\pmod N$.
\end{lemma}
\begin{proof}
For an odd prime $\ell$, the sets of squares and of $-1$
minus squares in $\mathbb F_\ell$ both have size
$(\ell+1)/2$, so they intersect.  Thus $-1$ is a sum of
two squares modulo $\ell$.  At least one of the two
coordinates is nonzero.  Adjusting that coordinate at each
step lifts the solution to $\ell^s$, since its derivative
$2a$ or $2b$ is a unit modulo $\ell$.
For powers of two, start with
$1^2+1^2+1^2+2^2=7\equiv-1\pmod8$.
If the sum is correct modulo $2^s$, $s\ge3$, replacing
the first, odd coordinate $a$ by $a+2^{s-1}$ changes
the sum by $2^s$ modulo $2^{s+1}$.  It therefore corrects
the next bit when necessary.  The same initial tuple
also works modulo $2$ and $4$.
Apply the Chinese remainder theorem to each of the four
coordinates.  The assertion for $N=1$ is immediate.
\end{proof}

\begin{theorem}\label{tc-gauss-eight}
If $q$ is nondegenerate, then
\[
                    \left(\sum_{x\in M}q(x)\right)^8=|M|^4.
 \tag{TC16}
\]
More generally, if its sum $S$ is nonzero, then $q$ is trivial
on its radical $R$ and
\[
                         S^8=|R|^4|M|^4>0.
 \tag{TC17}
\]
\end{theorem}
\begin{proof}
First suppose $R=0$, and set $N=|M|^2$.
Choose $a,b,c,d$ as in Lemma~\ref{tc-four-squares-congruence}
and write $s=a^2+b^2+c^2+d^2$.
The integer quaternion matrix
\[
 T=\begin{pmatrix}
 a&-b&-c&-d\\
 b&a&-d&c\\
 c&d&a&-b\\
 d&-c&b&a
 \end{pmatrix}
 \quad\hbox{satisfies}\quad
                  T^{\mathsf T}T=TT^{\mathsf T}=sI_4.
 \tag{TC18}
\]
It permutes $M^4$, because $s\equiv-1\pmod N$ and
$N$ annihilates $M$; its inverse on this group is
$-T^{\mathsf T}$.
For $Q(x_1,\ldots,x_4)=\prod_iq(x_i)$ the quadratic
identity (TC15), applied to each row, gives
\[
                         Q(Tx)=Q(x)^s=Q(x)^{-1}.
 \tag{TC19}
\]
Indeed the exponents of $q(x_i)$ are the squared column
lengths $s$, while the exponents of each $b(x_i,x_j)$
are the column inner products, all zero for $i\ne j$.
The last equality uses $q(x_i)^N=1$.
Summing (TC19) gives $S^4=\overline S^{\,4}$.
On the other hand
\[
 |S|^2
  =\sum_{t\in M}q(t)\sum_{y\in M}b(y,t)=|M|.
 \tag{TC20}
\]
Here we substituted $x=y+t$ in $\sum_{x,y}q(x)/q(y)$;
nondegeneracy and character orthogonality leave only $t=0$.
Multiplying $S^4=\overline S^{\,4}$ by $S^4$ proves (TC16).

For a possibly degenerate $q$ with $S\ne0$, translation
of its sum by $r\in R$ gives $S=q(r)S$.
Thus $q(r)=1$.  The function descends to a nondegenerate
quadratic function $\bar q$ on $M/R$, and
$S=|R|\sum_{M/R}\bar q$.
Equation (TC16) gives
$S^8=|R|^8|M/R|^4=|R|^4|M|^4$, as claimed.
\end{proof}

\subsection{The phase in a composition of three lattice averages}

\begin{theorem}\label{tc-triple}
For three commensurable self-dual lattices with a common sign
semicharacter there is a scalar $\mu$ such that
\[
                      I_{23}I_{12}=\mu I_{13},
                         \qquad \mu^8=1
 \tag{TC21}
\]
on every function fixed by $\Lambda_1$.
\end{theorem}
\begin{proof}
Choose representatives $S_2$ for
$\Lambda_2/(\Lambda_1\cap\Lambda_2)$ and $S_3$ for
$\Lambda_3/(\Lambda_2\cap\Lambda_3)$.
For $z\in\Lambda_3$, $y\in\Lambda_2$, $z+y\in\Lambda_1$
put
\[
 \psi(z,y)=\alpha(z)\alpha(y)\alpha(z+y)e^{-\pi iE(z,y)},
 \quad
 B=\sum_{\substack{z\in S_3,\ y\in S_2\\z+y\in\Lambda_1}}
                         \psi(z,y).
 \tag{TC22}
\]
The projection argument of Lemma~\ref{tc-unitary-intertwiner}
gives, for $F$ fixed by $\Lambda_1$,
\[
 I_{31}I_{23}I_{12}F
       =\frac{\sqrt{d_{13}}}{\sqrt{d_{12}}\sqrt{d_{23}}}\,BF.
 \tag{TC23}
\]
In detail, expand the two rightmost averages and apply
$I_{31}=\sqrt{d_{13}}P_1$ on their output.
The projection kills $W_zW_yF$ unless $z+y\in\Lambda_1$,
by character orthogonality.  In the remaining cases
$W_zW_yF=\psi(z,y)F$ by (TC4).
Applying the inverse $I_{13}$ proves (TC21) with
$\mu=\sqrt{d_{13}}B/(\sqrt{d_{12}}\sqrt{d_{23}})$.
The unitarity of all three operators, applied to a nonzero
theta function from Lemma~\ref{tc-covolume}, gives $|\mu|=1$.

It remains to determine its phase.  Define subgroups
\[
 \begin{split}
 P&=\{(z,y)\in\Lambda_3\times\Lambda_2:z+y\in\Lambda_1\},\\
 T&=\{(t,u-t):t\in\Lambda_2\cap\Lambda_3,\
                         u\in\Lambda_1\cap\Lambda_2\}.
 \end{split}
 \tag{TC24}
\]
The pairs in the sum (TC22) form representatives for $P/T$.
To see existence, first reduce $z$ modulo
$\Lambda_2\cap\Lambda_3$ to $z_0\in S_3$, moving the
discarded element into $y$, and then reduce the new $y$
modulo $\Lambda_1\cap\Lambda_2$ to $y_0\in S_2$.
The sum $z_0+y_0$ stays in $\Lambda_1$.
For uniqueness, membership of the difference in $T$ first
forces equality of the $S_3$ representatives and then
equality of the $S_2$ representatives.  In particular $P/T$
is finite.

Direct use of (TC2) gives
\[
 \begin{split}
 \psi(z+z',y+y')&=\psi(z,y)\psi(z',y')
                                      e^{2\pi iE(y,z')},\\
 \psi(nz,ny)&=\psi(z,y)^{n^2}.
 \end{split}
 \tag{TC25}
\]
For the first equation, the excess exponent after expanding
the three semicharacters is
$2\pi i\{E(z,z')+E(y,y')+E(y,z')\}$.
The first two terms are integers and disappear.
For the second, $\alpha(nx)=\alpha(x)^n
=\alpha(x)^{n^2}$ on each lattice, since its values are signs.
The polar factor in (TC25) is a bicharacter.
For $(t,u-t)\in T$, formula (TC2) on $\Lambda_2$ gives
$\psi(t,u-t)=1$, and $E(y,t)\in\mathbb Z$ gives trivial
polar factor against every $(z,y)\in P$.
Thus $\psi$ descends to a homogeneous quadratic function
on $P/T$.  Its sum is $B$.
Since $|\mu|=1$, this sum is nonzero.  By
Theorem~\ref{tc-gauss-eight}, $B^8$ is positive real.
The multiplier relating $\mu$ to $B$ is positive real,
so $\mu^8$ is positive real and has absolute value one.
It is therefore one.
\end{proof}

\subsection{The Cayley phase and the Gaussian trace}

Let $\delta$ be a complex linear $H$-isometry with no nonzero
fixed vector.  Assume that $\Lambda$ and $\delta\Lambda$
are commensurable self-dual lattices, that the common function
$\alpha$ is a sign semicharacter on both, and that
\[
                         \alpha(\delta x)=\alpha(x)
                                      \quad(x\in\Lambda).
 \tag{TC26}
\]
This last equality is essential; agreement on an intersection
alone would not assert it.
Set
\[
 A=1-\delta^{-1},\quad C=A^{-1}-\tfrac12,\quad
 \Gamma=\Lambda\cap\delta\Lambda,\quad d=[\Lambda:\Gamma],
 \quad
 q(x)=\alpha(x)e^{-\pi H(Cx,x)}\quad(x\in\Lambda).
 \tag{TC27}
\]

\begin{lemma}\label{tc-cayley-quadratic}
The function $q$ descends to a nondegenerate homogeneous
quadratic function on $\Lambda/A\Gamma$.  Moreover
\[
                 N:=|\Lambda/A\Gamma|=d|\det_{\mathbb C}A|^2.
 \tag{TC28}
\]
\end{lemma}
\begin{proof}
Unitarity gives
$A^*=-\delta A$ and $A^{-1}+(A^*)^{-1}=1$,
so $C^*=-C$.
Thus $H(Cx,x)$ is purely imaginary.  For $x,y\in\Lambda$,
expanding $q(x+y)$ and using $E(x,y)\in\mathbb Z$ gives
\[
 \frac{q(x+y)}{q(x)q(y)}
                   =e^{-2\pi iE(A^{-1}x,y)}.
 \tag{TC29}
\]
For example the exponent before reducing modulo $2\pi i\mathbb Z$
is $\pi iE(x,y)-2\pi iE(Cx,y)$, equal to the displayed
exponent plus $2\pi iE(x,y)$.
This is a bicharacter; its symmetry follows also from its
definition as a polar factor.  Homogeneity follows from
$\alpha(nx)=\alpha(x)^{n^2}$.

If $t\in\Gamma$, then $t,\delta^{-1}t\in\Lambda$ and
$At=t-\delta^{-1}t\in\Lambda$.
Also
$H(CAt,At)=-iE(t,\delta^{-1}t)$.
Equations (TC2) and (TC26) give
\[
 q(At)=\alpha(t)\alpha(\delta^{-1}t)=1.
 \tag{TC30}
\]
The polar factor against $At$ is one, because $t\in\Lambda$.
Thus $q$ descends modulo $A\Gamma$.
The radical of (TC29) on $\Lambda$ consists precisely of
the $x$ with $A^{-1}x\in\Lambda^*=\Lambda$.
For such an $x$, write $t=A^{-1}x$.
The equation $\delta^{-1}t=t-x\in\Lambda$ implies
$t\in\Gamma$, hence $x\in A\Gamma$.
This proves nondegeneracy on the quotient.
Finally $A$ changes real volume by $|\det_{\mathbb C}A|^2$.
Lemma~\ref{tc-covolume} and $\operatorname{vol}(V/\Gamma)=d$
therefore give (TC28).
\end{proof}

\begin{theorem}\label{tc-unitary-constant}
For every $F\in\mathcal T(\Lambda,\alpha)$,
\[
 F(\delta^{-1}z)=
       c(\delta,\Lambda)\,I_{\Lambda,\delta\Lambda}F(z),
 \quad
 c(\delta,\Lambda)=
 \frac{|\det A|}{\det A}\,
         \frac1{\sqrt N}\sum_{x\in\Lambda/A\Gamma}q(x).
 \tag{TC31}
\]
In particular
\[
                         c(\delta,\Lambda)^8=(\det\delta)^4.
 \tag{TC32}
\]
\end{theorem}
\begin{proof}
Both sides of the first equation belong to the same
one-dimensional theta space; the left side has the required
transformation law by (TC26).  Lemma~\ref{tc-unitary-intertwiner}
makes its right-hand operator nonzero.
There is therefore a scalar $c$, independent of $F$.
We compute it, which is the step needed for its arithmetic
absolute value.

Put $K_w(z)=e^{\pi H(z,w)}$ and consider the periodized kernel
\[
                       P(z,w)=\sum_{x\in\Lambda}W_xK_w(z).
 \tag{TC33}
\]
For $z,w$ in compact sets, each summand and every derivative
is bounded by a polynomial in $\|x\|$ times
$e^{-\pi\|x\|^2/2+C\|x\|}$.
The lattice sum is consequently normally convergent.
Reindexing under $W_\lambda$ shows that
$P(\,\cdot\,,w)\in\mathcal T(\Lambda,\alpha)$.

We need a diagonal integral.  For $y\in\delta\Lambda$,
expand (TC33) and use (TC4).  After multiplication by
$e^{-\pi H(z,z)}$, its $x$-summand is
\[
 \alpha(y)\alpha(x)e^{-\pi iE(y,x)}
       e^{-\pi H(x+y,x+y)/2}e^{2\pi iE(x+y,z)}.
 \tag{TC34}
\]
The last factor is a character of the real torus $V/\Gamma$.
It has integral zero unless $x+y=0$; nondegeneracy of $E$
ensures that a nonzero $x+y$ gives a nontrivial character.
If $y\in\Lambda$ the unique surviving term is $x=-y$,
whose coefficient is one.  Otherwise no term survives.
Termwise integration is justified by the preceding Gaussian
bound on a compact fundamental parallelepiped.  Thus
\[
 \int_{V/\Gamma}(W_yP(\,\cdot\,,z))(z)
                        e^{-\pi H(z,z)}\,dz
     =\begin{cases}d,&y\in\Lambda,\\0,&y\notin\Lambda.\end{cases}
 \tag{TC35}
\]
The case $y=0$ shows that the kernel is not identically zero.
Averaging (TC35) over representatives of $\delta\Lambda/\Gamma$
gives
\[
 \int_{V/\Gamma}
   (I_{\Lambda,\delta\Lambda}P(\,\cdot\,,z))(z)
                          e^{-\pi H(z,z)}\,dz=\sqrt d.
 \tag{TC36}
\]

Set $g(z)=e^{-\pi H(Az,z)}$.
A direct expansion, using $C^*=-C$, gives the exact identity
\[
 \alpha(x)(U_xK_z)(\delta^{-1}z)e^{-\pi H(z,z)}
                  =q(x)\,g(z-A^{-1}x).
 \tag{TC37}
\]
For clarity, the linear terms agree because
$H(Az,A^{-1}x)=-H(\delta^{-1}z,x)$.
The constant terms agree because
$H(A^{-1}x,x)+H(x,A^{-1}x)=H(x,x)$.
Also
\[
                 |g(z)|=e^{-\pi H(Az,Az)/2}.
 \tag{TC38}
\]
This follows from $A+A^*=A^*A$ and is an integrable
positive Gaussian, since $A$ is invertible.
Decompose $\Lambda$ into the $N$ cosets of $A\Gamma$.
By Lemma~\ref{tc-cayley-quadratic}, the phase $q$ is constant
on each coset.  The integral of
$\sum_{\gamma\in\Gamma}g(z-A^{-1}x-\gamma)$ over
a fundamental domain is $\int_Vg$.
Absolute convergence follows from (TC38); hence both the
decomposition and this unfolding are justified.
We obtain
\[
 \int_{V/\Gamma}P(\delta^{-1}z,z)e^{-\pi H(z,z)}\,dz
       =\left(\sum_{\Lambda/A\Gamma}q(x)\right)\int_Vg(z)\,dz.
 \tag{TC39}
\]

Diagonalize the unitary map $\delta$ in an orthonormal basis.
The eigenvalues of $A$ are $\lambda_j=1-e^{-i\theta_j}$,
with $\operatorname{Re}\lambda_j>0$, since $\delta$ has no
fixed vectors.  Polar coordinates give, for such a $\lambda$,
\[
 \int_{\mathbb C}e^{-\pi\lambda|w|^2}\,dw
        =2\pi\int_0^\infty e^{-\pi\lambda t^2}t\,dt
        =\lambda^{-1}.
\]
Absolute convergence permits multiplication of these integrals,
so $\int_Vg=(\det A)^{-1}$.
Comparing (TC36) and (TC39) in the scalar identity
$P(\delta^{-1}z,w)=cI_{\Lambda,\delta\Lambda}P(\,\cdot\,,w)(z)$
now gives
\[
                c=\frac{\sum q(x)}{\sqrt d\,\det A}.
\]
Using $\sqrt N=\sqrt d\,|\det A|$ proves (TC31).

The normalized Gauss sum in (TC31) has eighth power one
by Theorem~\ref{tc-gauss-eight}.
Taking determinants in $A^*=-\delta A$ gives
\[
             \overline{\det A}=(-1)^r\det(\delta)\det A.
 \tag{TC40}
\]
Consequently $(|\det A|/\det A)^2=(-1)^r\det\delta$,
whose fourth power is $(\det\delta)^4$.
This proves (TC32).
\end{proof}

\subsection{Application to the actual fourth-power normalization}

Let $k=\mathbb Q(\sqrt{-D})$ with $D>0$, $D\equiv7\pmod8$,
and $\mathcal O_k=\mathbb Z[\tau]$,
$\tau=(1+i\sqrt D)/2$.
The same argument applies whenever the specified order is
$\mathbb Z[\tau]$; here $\mathcal O_k$ denotes the maximal
order used in the arithmetic construction.
Suppose $\Lambda\subset V$ is a full lattice stable under
$\mathcal O_k$ and self-dual for $E$.
Define on the whole complex space
\[
              \alpha_H(x)=
                    \exp\!\left(\pi i\frac{\sqrt D}{2}H(x,x)\right).
 \tag{TC41}
\]

\begin{lemma}\label{tc-norm-semicharacter}
The function (TC41) is a sign semicharacter on $\Lambda$.
Every isometry preserves this function, and it is multiplicative
under orthogonal direct sums.
On the standard lattice $\mathcal O_k^m$ with
$H_0(z,w)=(2/\sqrt D)\sum_jz_j\overline w_j$, it is
$(-1)^{\sum_jN_{k/\mathbb Q}(x_j)}$.
\end{lemma}
\begin{proof}
Put $Q(x)=(\sqrt D/2)H(x,x)$.
Since $\tau x\in\Lambda$,
$Q(x)=E(\tau x,x)\in\mathbb Z$.
For $x,y\in\Lambda$,
\[
 Q(x+y)-Q(x)-Q(y)+E(x,y)=2E(\tau x,y)\in2\mathbb Z.
 \tag{TC42}
\]
Subtracting $2E(x,y)\in2\mathbb Z$ gives (TC2).
The remaining assertions follow immediately from preservation
and additivity of the squared Hermitian norm.
\end{proof}

We also make explicit the nonprincipal elliptic base allowed
in the submitted construction.
Let
\[
 I=\mathbb Za+\mathbb Z\tau_b,\qquad
 \tau_b=\frac{b+i\sqrt D}{2},\qquad
 4ac=b^2+D,\qquad a>0\text{ odd},
 \tag{TC43}
\]
be an integral $\mathcal O_k$-ideal, and let $L=\lambda I$,
$\lambda\in\mathbb C^\times$.
Its Hermitian form is
$H_L(z,w)=2z\overline w/(a|\lambda|^2\sqrt D)$.

\begin{lemma}\label{tc-ideal-character}
The normalized translate by $\lambda a/2$ of the reduced odd
theta function of $L$ has semicharacter $\alpha_{H_L}$.
The lattice $L$ is self-dual for $\operatorname{Im}H_L$.
\end{lemma}
\begin{proof}
In the basis $\lambda a,\lambda\tau_b$, the alternating
pairing has entries $E(\lambda a,\lambda\tau_b)=-1$ and
$E(\lambda a,\lambda a)=E(\lambda\tau_b,\lambda\tau_b)=0$.
It is therefore unimodular.
For $x=\lambda(am+n\tau_b)$ the reduced odd theta
semicharacter is $(-1)^{m+n+mn}$.
Translation by $\lambda a/2$ multiplies it by
$e^{2\pi iE(x,\lambda a/2)}=(-1)^n$, yielding
$(-1)^{m+mn}$; the opposite sign in this exponential
gives the same value.
On the other hand
\[
 \frac{\sqrt D}{2}H_L(x,x)
       =\frac{N(am+n\tau_b)}a
       =am^2+bmn+cn^2.
 \tag{TC44}
\]
The relation in (TC43) makes $b$ odd, and then
$D\equiv7\pmod8$ and oddness of $a$ make $c$ even.
Thus (TC44) has parity $m+mn$, as required.
The translating half-period does not belong to $L$, so the
value used to normalize the theta function at zero is nonzero.
\end{proof}

Now put $\Lambda_0=L^{4r}$ with the product form $H_0$,
and let $\Theta_0$ be the product of these normalized
one-dimensional theta functions.
The case $I=\mathcal O_k$ and $\lambda=1$ is the base of
Section~\ref{ct-section}.  Lemma~\ref{tc-ideal-character}
handles the other ideal classes without changing the
analytic argument.
Let $g_0,g_1:V^4\to\mathbb C^{4r}$ be isometries from
$H^{\oplus4}$ to $H_0$, and put $M_i=g_i(\Lambda^4)$.
Assume that $M_0,M_1,\Lambda_0$ are commensurable and
$\det(g_1g_0^{-1})$ is algebraic.  These are satisfied
when the maps are obtained from isometries over $k$ in
the common scaled $k$-coordinates of the construction.
Use the same norm semicharacter $\alpha_{H_0}$ on all
three lattices.  They are self-dual by the isometries and
Lemma~\ref{tc-ideal-character}.
Define, as in Section~\ref{lm-section},
\[
          K_{g_i}(z)=(I_{\Lambda_0,M_i}\Theta_0)(g_i z).
 \tag{TC45}
\]

\begin{theorem}\label{tc-kernel-ratio}
Suppose $g_1-g_0$ is invertible.  Then there is a nonzero
algebraic number $c$ such that
\[
       K_{g_0}=cK_{g_1},\qquad
       c^8=\det(g_1g_0^{-1})^4.
 \tag{TC46}
\]
Thus the analytic identity required in the fourth tensor
criterion of Section~\ref{lm-section} holds with its exact
determinant exponent.
\end{theorem}
\begin{proof}
Put $\delta=g_1g_0^{-1}$ and $F=I_{\Lambda_0,M_0}\Theta_0$.
This is a nonzero theta function by
Lemma~\ref{tc-unitary-intertwiner}.
The hypothesis makes $\delta-1$ invertible.
The norm semicharacter is invariant under $\delta$ by
Lemma~\ref{tc-norm-semicharacter}, including the equality
(TC26).  Apply Theorem~\ref{tc-unitary-constant} on $M_0$
and then Theorem~\ref{tc-triple}:
\[
 F(\delta^{-1}w)
    =c_0 I_{M_0,M_1}I_{\Lambda_0,M_0}\Theta_0(w)
    =c_0\mu I_{\Lambda_0,M_1}\Theta_0(w),
 \quad c_0^8=(\det\delta)^4,\quad\mu^8=1.
\]
Evaluate at $w=g_1z$.  This gives (TC46) with $c=c_0\mu$.
Since $\det\delta$ is algebraic, that equation also proves
that $c$ is algebraic.  No algebraicity of the separate
complex Gaussian factors has been presumed.
\end{proof}

\begin{corollary}\label{tc-normalized-unit}
Use the reference isometry $\gamma_0=g_0$, and form
$\kappa$ as in (LM9)--(LM10), with the locally standard
maps carrying the given lattice onto the fixed target
$\Lambda_0$ in scaled $k$-coordinates.
If $p$ splits in $k$ and
$g_1(\Lambda^4\otimes\mathbb Z_p)
=\Lambda_0\otimes\mathbb Z_p$ in those coordinates, then,
for the same
embedding into $\mathbb C_p$ as in that theorem,
\[
                         |\iota(\kappa c)|_p=1.
 \tag{TC47}
\]
Moreover a nonzero theta function $F_\Lambda$ with
$F_\Lambda^{\otimes4}=\kappa K_{g_0}$ exists and is unique
up to a fourth root of unity.
\end{corollary}
\begin{proof}
The proof of Theorem~\ref{lm-kappa} applies to any fixed
target $\mathcal O_k$-lattice.  Indeed, two local maps
onto that target differ by an automorphism of its completed
lattice; its determinant has valuation zero at both primes.
The finite-support argument uses commensurability with the
fixed target in place of commensurability with
$\mathcal O_k^{4r}$, and every subsequent valuation
calculation is unchanged.
Equation (TC46) and Corollary~\ref{lm-change-scalar}
therefore give (TC47).  For existence, choose any nonzero element $\theta$
of $\mathcal T(\Lambda,\alpha_H)$ by
Lemma~\ref{tc-covolume}.
Both $\theta^{\otimes4}$ and $K_{g_0}$ are nonzero elements
of the one-dimensional theta space for $\Lambda^4$.
Thus $\kappa K_{g_0}=a\theta^{\otimes4}$ for $a\ne0$.
Choose a fourth root of $a$ and multiply $\theta$ by it.
Any two choices, and any two theta functions satisfying
the displayed normalization, differ by a fourth root of
unity, again by the one-dimensionality of the theta space.
\end{proof}

The determinant exponent in (TC46), the positive index
square roots in (TC10), and the invariant sign semicharacter
in (TC26) have all been retained.  The complex absolute
value of a theta normalization alone would not establish
(TC47).  The calculation above establishes the stronger
algebraic equality that supplies the finite-place bound.
It closes the analytic change-of-isometry input; identifying
the particular arithmetic column and its support is a
separate subsequent step.


